\documentclass{beamer}
%\usepackage{split}

\usepackage{amsmath}
\usepackage{amsfonts}
\usepackage{braket}
\usepackage{subcaption}
\usepackage{tikz}
\usepackage{xcolor}
\usepackage{svg}
\usetikzlibrary{positioning,shapes,arrows.meta}
%\usetikzlibrary{positioning,fit,backgrounds}
\usepackage{quantikz}
\usepackage[linesnumbered]{algorithm2e}
\usepackage{hyperref}
%\usepackage{hyperref,xcolor}
%\usepackage[ocgcolorlinks]{ocgx2}
\usepackage{cleveref}
\usepackage[backend=biber,style=alphabetic,sorting=ynt]{biblatex}
\addbibresource{./sample.bib}


\input{newcommands}


\newcommand*{\QACze}{ \mathbf{QAC}_{0} }
\newcommand*{\QNCzef}{ \mathbf{QNC}_{0,f} }
\newcommand*{\QNCze}{ \mathbf{QNC}_{0} }
\newcommand*{\QNCon}{ \mathbf{QNC}_{1} }
\newcommand*{\NCon}{ \mathbf{NC}_{1} }
\newcommand*{\noiseQNCon}{ noisy-$\QNCon$ }
\newcommand*{\QNC}{ \mathbf{QNC} }
\newcommand*{\QNCG}{ \mathbf{QNC_G} }
\newcommand*{\NC}{\mathbf{NC}}
\newcommand*{\QNCiG}{\mathbf{QNC_{G,i}}}






\usetikzlibrary{patterns, shapes.arrows}
\usepackage{tkz-graph}
\usetikzlibrary{decorations.pathmorphing}
\usepackage{tkz-berge}

\usetikzlibrary{intersections}

\usepackage{cancel}
\begin{document} 

\newcommand*{\Tr}{\textbf{Tr }}


\begin{frame}
  \title{Topological Memory.}
    \author{David Ponarovsky}
    \date{\today}
    \titlepage
\end{frame}


\begin{frame}

\frametitle{Introduction}
\begin{block}{Today:}
\begin{itemize}
  \item Proof. 
    \begin{enumerate}
      \item Walktrough the math 
      \item A 
      \item B
    \end{enumerate}
  \item New experimental results.  
  \item Answering the code teleportation.  
\end{itemize}
\end{block}
\end{frame}

% ------------------------------------------------ $ 


\begin{frame}
  \frametitle{Our Code}

  \begin{block}{Code}
    In this session, we consider LDPC codes, where each bit adjoins to $\Delta_{1}$ checks, and each check adjoins to $\Delta_{2}$ bits. We denote by $\Delta = \Delta_{1}\Delta_{2}$.
  \end{block}


\end{frame}

\begin{frame}
  \frametitle{Local Decoder} 
  
  \begin{block}{definition}
    We say that $D$ is a local decoder if when applied against $2p(1-p)$-depolarized channel: 
    \begin{itemize}
      \item  There exists function $M : (0,1) \rightarrow (0,1) $ such that the probability of each bit being flipped is $M(2p(1-p)) < p$.
      \item There is a constant $l$ such that for any bit, there are only $(\Delta)^{l}$ bits on which the event of it being flipped depends.
    \end{itemize}
For example, if we take $k$ large enough and consider the $(n,k)$-Toric, then both the swift rule-and the flip upon majority are local decoders.
  \end{block}
\end{frame}

\begin{frame}
  \frametitle{Storing Logical State}

  \begin{block}{Logical State}
    For a random variable $X$ we denote by: 
    \begin{enumerate}
      \item $\exppm{X}{ \sim p  }$ - its expectation when measured after a single application of the $p$-depolarizing channel.
      \item $\exppm{X}{ \sim q^{(t)}  }$ - its expectation when measured after $t$ rounds, each initiated by a $p$-depolarized channel followed by an application of local decoder $D$.
    \end{enumerate}
  \end{block}
\end{frame}


\begin{frame}  
  \frametitle{The Theorem} 
  \begin{block}{Theorem} Let $f$ be any function with bounded value, denote by $|f|$ it's supremum, Then:
    \begin{equation*}
      \exppm{f}{\sim q^{(t)} }  \le \exppm{ f }{ \sim p  } + t|f| e^{- \Theta( \sqrt{n}) }
    \end{equation*}
\end{block}

\begin{block}{Lemma} There exists a constant $c >0 $, such:
    \begin{equation*}
        \exppm{ f \circ D }{\sim 2p(1-p) } \le \exppm{ f \circ D }{\sim p } + |f| e^{- c \sqrt{n} } 
    \end{equation*}
\end{block}
\end{frame}

\begin{frame}
  \frametitle{Theorem Proof}
  By induction on $t$.
  \begin{align*}
    \exppm{f}{\sim q^{(t)} } &=  \exppm{ \exppm{ f \circ D }{\sim p}  }{\sim q^{(t-1)} } \\
      & \le   \exppm{ \exppm{ f \circ D }{\sim p}  }{\sim p } + (t-1)| \exppm{f \circ D}{p}   | e^{- c \sqrt{n} } \\
      & \le   \exppm{ f \circ D }{\sim 2p(1-p) } + (t-1)|f| e^{- c \sqrt{n} } \\
      & \le \exppm{ f }{ \sim p  } + t|f| e^{- c \sqrt{n} }
  \end{align*}
  When in the last transition we used the lemma. 
\end{frame}


\begin{frame}

\end{frame}


\end{document}
